\section{Introduction}\label{sec:intro}

The $n$-queens problem asks how many ways there are to place $n$ queens on an $n\times n$ board
so that no two attack each other. We denote this number by $Q(n)$. The survey of Bell and
Stevens~\cite{BellStevens} covers the history of the problem and its variants. In a solution, each row
contains one queen. Recording its column as we read the rows from top to bottom gives a
permutation, since no column occurs twice. Figure~\ref{fig:permutaciones} shows why only some of the $n!$
possible
permutations are solutions.

\begin{figure}[htbp]
\centering
\begin{tikzpicture}[x=5.6mm,y=5.6mm]
  \begin{scope}
    \bcuadricula{4}
    \foreach \r/\c in {1/2,2/4,3/1,4/3} {\breina{4}{\r}{\c}}
    \node[font=\small,anchor=north] at (2,-.35) {(a) Solution};
  \end{scope}
  \begin{scope}[shift={(9,0)}]
    \bcuadricula{4}
    \draw[btinta,densely dashed,line width=1.1pt] (.5,3.5)--(3.5,.5);
    \foreach \r/\c in {1/1,2/2,3/3,4/4} {\breina{4}{\r}{\c}}
    \node[font=\small,anchor=north] at (2,-.35) {(b) Shared diagonal};
  \end{scope}
\end{tikzpicture}
\caption{Two permutations on the same board. Rows are numbered from top to bottom and columns
from left to right, starting at one. In (a) the column sequence is $(2,4,1,3)$: the four
queens also lie on distinct diagonals in each direction. In (b) it is $(1,2,3,4)$:
each row and each column still has one queen, but the diagonal shown contains all four.}
\label{fig:permutaciones}
\end{figure}


By how much do the diagonals reduce these $n!$ possibilities as the board grows?
Simkin~\cite[Theorem~1.1]{Simkin} determined the leading order by a variational principle:
there is a constant $\alpha$ such that
\[
 Q(n)=\bigl((1+o(1))\,n\,e^{-\alpha}\bigr)^n .
\]
Nobel, Agrawal and Boyd~\cite{NAB} bounded $\alpha$ in $[1.944000752,\,1.944001082]$.
Taking logarithms and applying Stirling's formula, Simkin's result gives
\[
 \log Q(n)-\log n!=n(1-\alpha)+o(n).
\]
The remainder $o(n)$ is undetermined. The question motivating this paper concerns the next term:
after subtracting the linear term, does the difference converge, and if so, to what constant?

We construct a deterministic quantity $Q_{\rm pred}(n)$, which we call the predictor and whose constant
term we can calculate. The construction assigns a weight to each square so that each row and
each column sums to one and no diagonal receives more than one. Among the admissible weights
of the model, we choose those that maximize the Bethe entropy (Section~\ref{sec:predictor}). This entropy
contributes the first term of the predictor. The other three terms in Definition~\ref{def:Qpred} are a
local term depending on the square weights, a pseudodeterminant that also depends on how the
lines intersect, and a normalization. The definition requires the maximizer to be interior,
which holds for all sufficiently large $n$ (Section~\ref{sec:P02}).

Passing from the board to the unit square brings in the profile of mass one $\nu_1$ from~\cite{Perfil}.
We write $s_1$ for the optimal value at mass one, attained by that profile, and set $\Phi=s_1+1$.
Section~\ref{sec:preliminares} deduces that $\Phi=1-\alpha$, where $\alpha$ is Simkin's constant;
no proof of the theorems uses this equality. The integral $I_2=\int\nu_1^2$ is taken over the unit square.
Since $\nu_1$ has total mass one,
$I_2\ge1$, with equality only for the constant density: $I_2$ measures how far the profile is from that
density.

\begin{teorema}[Constant term of the predictor]\label{thm:intro-constante}
As $n\to\infty$,
\[
\begin{aligned}
 \log Q_{\rm pred}(n)-\log n!-n\Phi &\longrightarrow C_{\rm pred},\\
 C_{\rm pred} &=\tfrac12\log2-\tfrac12-\tfrac32I_2-\tfrac12\log{\det}_2(I+J),
\end{aligned}
\]
where $J$ is the Hilbert--Schmidt operator constructed from $\nu_1$ in Section~\ref{sec:D03} and
$I+J$ is positive definite.
\end{teorema}

The proof follows the definition term by term. First we approximate the discrete maximizer by the
continuous profile with error $O(1/n)$ (Section~\ref{sec:perfil}). We then identify the contributions of the
entropy and the pseudodeterminant separately and add them to the local term
(Section~\ref{sec:constante}). The theorem is proved in Section~\ref{sec:D05} (Theorem~\ref{thm:constante}).
Of the continuous profile, the proof uses only the properties taken from~\cite{Perfil} that
Section~\ref{sec:preliminares} collects: existence and uniqueness of $\nu_1$, the
factorization~\eqref{eq:factorizacion}, its bounds and a Lipschitz constant. Here $\det_2$ is the
regularized determinant of Hilbert--Schmidt operators, defined in Section~\ref{sec:D04}
(see~\cite[Ch.~9]{Simon}).

The second result answers a question that can already be posed on a small board.
A partial placement leaves a residual board $B$. We choose an empty row $r$ and write
$B+c$ for the residual board left after occupying a legal square $c$ in that row.
If $Q(B)$ counts the completions of $B$, then
\[
 Q(B)=\sum_{c\in r} Q(B+c).
\]
Indeed, each completion uses exactly one of these squares. A legal choice that has no completion
contributes zero to the sum. Figure~\ref{fig:residual} shows this partition on a $5\times5$ board with a
queen in
row~1 and column~2: two legal squares remain in the second row, in columns 4 and~5, and each gives one
summand.

\begin{figure}[htbp]
\centering
\begin{tikzpicture}[x=5.6mm,y=5.6mm]
 \begin{scope}
  \bcuadricula{5}
  \foreach \r in {1,...,5} {\foreach \c in {1,...,5} {
    \pgfmathtruncatemacro{\ataque}{(\r==1)||(\c==2)||(abs(\r-1)==abs(\c-2))}
    \ifnum\ataque=1\relax\bmarca{5}{\r}{\c}\fi
  }}
  \breina{5}{1}{2}
  \foreach \c in {4,5}{\draw[btinta,line width=1.1pt] (\c-.5,3.5) circle[radius=.31];}
  \node[anchor=north,font=\small] at (2.5,-.32) {(a) Parent $B$};
 \end{scope}
 \draw[-{Stealth[length=2mm]},btinta] (5.55,2.5)--(7.35,2.5);
 \node[font=\footnotesize,above] at (6.45,2.65) {$(2,4)$};
 \begin{scope}[shift={(8.5,0)}]
  \bcuadricula{5}
  \foreach \r in {1,...,5}{\foreach \c in {1,...,5}{
    \pgfmathtruncatemacro{\ataque}{(\r==1)||(\c==2)||(abs(\r-1)==abs(\c-2))||(\r==2)||(\c==4)||%
(abs(\r-2)==abs(\c-4))}
    \ifnum\ataque=1\relax\bmarca{5}{\r}{\c}\fi
  }}
  \breina{5}{1}{2}\breina{5}{2}{4}
  \node[anchor=north,font=\small] at (2.5,-.32) {(b) Child $B+(2,4)$};
 \end{scope}
\end{tikzpicture}
\caption{The exact partition of the count on a small board. The second row of the board on the
left has two legal choices, marked with rings. The arrow follows one of them; the other,
$(2,5)$, gives the second summand of the count. Dots mark occupied squares or squares under attack,
and crowns denote queens already placed. The available squares retain their coordinates.}
\label{fig:residual}
\end{figure}


The predictor assigns to $B$ the value $e^{L(B)}$ (Section~\ref{sec:B3T01}).
We divide the sum of the predictions for the children by the prediction for the parent:
\[
 \rho(B)=\frac{\sum_{c\in r}e^{L(B+c)}}{e^{L(B)}}
        =\sum_{c\in r}e^{L(B+c)-L(B)} .
\]
If the predictor satisfied exactly the same identity as the count, this ratio would be one.
The ratio also depends on the chosen row $r$, which we do not indicate in the notation.
The following theorem bounds the defect on a logarithmic scale and specifies the residual boards for which
the bound holds.

Its statement uses notions made precise in Sections~\ref{sec:insercion} and~\ref{sec:B5T01}.
A residual board with $k$ empty rows also has $k$ empty columns and is missing $k$ queens;
in \cite{Completabilidad} this number is denoted by $h$.
The \emph{Bethe point} of $B$ is the vector of weights $p=(p_e)$, one for each available square
$e$, at which $L(B)$ is evaluated: a stationary point of the Bethe entropy of the residual board
under the row and column constraints. The \emph{mass} of a diagonal in either family is the sum
of the weights of its available squares; its \emph{vacancy} is one minus its mass.
The \emph{row--column Laplacian} is the Laplacian of the graph whose vertices are the empty rows
and columns and in which each available square $e$ is an edge of weight $p_e$ between its row and
its column. Its \emph{gap} is its second smallest eigenvalue, counted with multiplicity.

\begin{teorema}[Second-order consistency]\label{thm:intro-consistencia}
Let $C,\eta,\gamma>0$ and let $B$ be a residual board with $k$ empty rows whose Bethe point $p$ is
stationary and interior, with row and column masses equal to one, $p_e\le C/k$, diagonal vacancies
at least $\eta$, and a gap of at least $\gamma$ for the row--column Laplacian weighted by $p$.
\begin{enumerate}[(a)]
\item For every $\delta_0>0$ there exist $k_0$ and $C_*$, depending only on $\delta_0,C,\eta,\gamma$, such
that if $k\ge\max(k_0,\delta_0n)$, then $|\log\rho(B)|\le C_*/k^2$ for every empty row $r$
(Theorem~\ref{thm:macro} and Corollary~\ref{cor:rho}).
\item There exist $\varepsilon_0$, $k_0$ and $C_*$, independent of $n/k$, such that the same conclusion
holds for $k\ge k_0$ without the condition $k\ge\delta_0n$ if the mass of each diagonal is at most
$\varepsilon_0$ (Theorem~\ref{thm:meso}).
\end{enumerate}
\end{teorema}

Both parts include all legal children of the chosen row. For each child we obtain a unique
stationary point on the small branch, that is, with all its weights less than $1/4$. The first
part requires a fixed fraction of the original rows to remain. The second dispenses with
this condition when the mass of each diagonal is small. In both cases the hypotheses concern
the residual board. The theorem is deterministic: it does not assert that the residual boards
arising when queens are placed at random satisfy these hypotheses.
The class in part~(a) is not empty: for large $n$ it contains the full board
(Corollary~\ref{cor:completo}). For the class in part~(b) we exhibit no residual of queens
with $k=o(n)$ (Section~\ref{sec:B5T05}).

To transfer the first theorem to the exact count, it remains to control a difference:
\[
 R_Q(n)=\log Q(n)-\log Q_{\rm pred}(n).
\]
If $R_Q(n)\to0$, we obtain
\[
 \log Q(n)=\log n!+n\Phi+C_{\rm pred}+o(1).
\]
We do not prove that $R_Q(n)\to0$. Section~\ref{sec:discusion} explains what additional estimates on the
trajectories
and on residual boards with few empty rows this step would require.

Numerically, $C_{\rm pred}\approx0.9401$: over the entire interval for $\alpha$ cited above, fits of the
predictor starting at $n=40$ or $60$ give values between $0.9400$ and $0.9402$, and a direct evaluation of
the formula gives $0.94014$ (Section~\ref{sec:numerico}); none of these figures is certified.
With the reference value $\Phi_0$ from that section, the exact counts for $16\le n\le27$ give values of
$\log Q(n)-\log n!-n\Phi_0$ between $0.913$ and $0.941$.
